%%%%%%%%%%%%%%%%%%%%%%%%%%%%%%%%%%%%%%%%%%%%%%%%%%%%%%%%%%%%%%%%%
% Methodology appendix. Static prose; live figures via \ilsa... macros.
%%%%%%%%%%%%%%%%%%%%%%%%%%%%%%%%%%%%%%%%%%%%%%%%%%%%%%%%%%%%%%%%%
\chapter{Methodology}
\label{app:methodology}

This appendix documents the data sources, the consolidation and classification
model, the statistical estimators, and the forward population model behind every
figure in this report. The toolchain treats the report not as a document that is
typed, but as a \emph{reproducible artifact built deterministically from primary
financial records}. The same inputs regenerate byte-comparable outputs on every
run, and no number in the narrative is entered by hand: each is computed by code
from the consolidated ledger or the original invoices and injected into the prose
as a \LaTeX{} macro. An auditor can therefore trace any figure backward ---
narrative $\rightarrow$ generated macro $\rightarrow$ generating function
$\rightarrow$ ledger or invoice row $\rightarrow$ original bank, PayPal, or DPIL
record.

\section{Design principles}

Three principles govern the design.

\begin{itemize}
  \item \textbf{Transparency.} All analysis reads one consolidated ledger
        (\texttt{ILSA\_full\_ledger.xlsx}) and the original Dollywood invoice
        PDFs. Transaction categorization lives in readable code, not in private
        spreadsheet judgment, so any reviewer can inspect and challenge the rules.
        Internal PayPal\,$\leftrightarrow$\,bank sweeps are detected and netted out
        so each gift is counted once.
  \item \textbf{Rigor.} Income is split into \emph{recurring} giving
        (forecastable) and \emph{episodic} grants and major gifts (shown discretely,
        never trended). Program cost is reported on an \emph{accrual} basis (the
        invoiced amount) and reconciled to \emph{cash} (bank payments). Every
        forward assumption is a named constant, and the projections carry an
        explicit $\pm$\ilsaScenBand{}/year uncertainty band.
  \item \textbf{Reproducibility.} A single shared model feeds the report, the
        figures, the tables, and the planning workbook, so they are mathematically
        incapable of disagreeing. The pipeline is parameterized by reporting period
        and is replayable end to end.
\end{itemize}

\section{Repository organization and build pipeline}

The project separates \emph{primary records}, a \emph{shared model}, and three
\emph{deliverables}. The shared library is never run directly; both the report
and the workbook import it, which is what guarantees they reconcile.

\begin{table}[h]
\centering
\small
\begin{tabular}{|l|p{9.2cm}|}
\hline
\textbf{Component} & \textbf{Role} \\
\hline
\texttt{ILSA\_full\_ledger.xlsx} & Single source of truth: bank (TCB) and PayPal exports, hand-merged. \\
\texttt{invoices/*.pdf} & Primary-source Dollywood invoices. \\
\texttt{parse\_invoices.py} & Parses each invoice PDF into \texttt{invoices/summary.csv}. \\
\texttt{ilsa\_ledger.py} & \textbf{Shared model.} Loads and consolidates the ledger, nets internal transfers, builds the monthly table, the invoice (accrual) model, and the projection. \\
\texttt{treasurer.py} & Builds the report: figures, CSV summaries, and \texttt{ILSA-commands.tex}. \\
\texttt{cohort\_explore.py} & Age-structured cohort model; generates the Section~2.3 figures and macros. \\
\texttt{running\_budget.py} & Builds the planning workbook \texttt{ILSA\_running\_budget.xlsx}. \\
\hline
\end{tabular}
\caption{Repository components and their roles.}
\end{table}

\noindent\textbf{Data flow.} Raw exports are merged into the consolidated ledger
and invoices are parsed into \texttt{summary.csv}. The shared library consolidates
both into the monthly model. \texttt{treasurer.py} emits the figures and a file of
\verb|\newcommand| definitions (\texttt{ILSA-commands.tex}); the prose files
(\texttt{FiscalReport.tex}, \texttt{2Report.tex}) reference only those macros, so
the wording stays editable while the numbers stay generated. The document is
compiled twice with \texttt{pdflatex} so the \texttt{lastpage} total page count
resolves.

\section{Data sources and consolidation}

Each ledger row is normalized to a common schema: posting \emph{date}, \emph{source}
(TCB or PayPal), \emph{description}, \emph{category}, signed \emph{amount} (net of
fees), and \emph{gross} (pre-fee). Bank amounts are signed by their credit/debit
flag; PayPal amounts use the \texttt{Net} field (fees included) with \texttt{Gross}
retained for donor-intent reporting. Rows are sorted by date and stamped with a
monthly period.

Invoices are parsed line by line. Each monthly invoice lists, per age group
(Group~1--6) and language, the book quantity and dollar amount, plus three
non-group items: welcome books (\texttt{LETC}), graduation books (\texttt{GRAD}),
and a per-piece mailing charge. The parser is robust to two invoice layouts,
recovers the printed total where present (falling back to the sum of line items),
and \emph{de-duplicates} re-exported invoices by their DPIL invoice number ---
keeping the most complete copy --- so a re-export can never double-count.
Promotional invoices (e.g.\ lapel pins) are detected and excluded from program
totals.

\section{Transaction classification}

Categorization is rule-based and auditable. Bank descriptions map to categories by
keyword: \texttt{DOLLYWOOD} $\rightarrow$ program cost; \texttt{BONTERRA}
$\rightarrow$ recurring giving; \texttt{DDA}/\texttt{DEPOSIT} $\rightarrow$ grant or
major gift; \texttt{FEE}/\texttt{SERVICE CHARGE} $\rightarrow$ bank fee;
\texttt{CHARGEBACK} $\rightarrow$ income reversal; \texttt{PAYPAL}/\texttt{INST XFER}
$\rightarrow$ internal transfer. PayPal rows are recurring giving unless they are
withdrawals, bank-funded deposits, or the pending bank-link wash, which are
internal. Internal transfers are excluded from both income and expense, so a
donation swept from PayPal to the bank is counted once (when received), never again
when the funds move.

\section{The invoice line-item model and enrollment conservation}

Reading the invoices as an age-structured population requires understanding how a
child appears on a given month's invoice. Empirically, the three recipient sets are
\emph{disjoint}: a continuing child receives a group book, a newly enrolled child
receives only a welcome book that month (entering an age group the following month),
and a graduating child receives a graduation book as they exit. Consequently the
mailing quantity equals the children served,
\begin{equation}
\text{kids}(M) \;=\; \underbrace{\textstyle\sum_{g} G_g(M)}_{\text{group books}}
              \;+\; \text{LETC}(M) \;+\; \text{GRAD}(M)
              \;=\; \text{mailing pieces}(M),
\end{equation}
and the group roll obeys a one-month conservation law,
\begin{equation}
\sum_g G_g(M) \;=\; \sum_g G_g(M-1) \;+\; \text{LETC}(M-1) \;-\; \text{GRAD}(M).
\label{eq:conservation}
\end{equation}
Equation~\eqref{eq:conservation} holds across the observed period with residuals of
a few children per month, which both validates the enrollment rule and confirms that
\texttt{LETC} is a clean measure of monthly inflow and \texttt{GRAD} of outflow.

\section{Cash versus accrual program cost}

The Dollywood Foundation \emph{invoices} a month's books (the cost incurred,
\emph{accrual}); the bank \emph{payment} clears about a month later (\emph{cash}).
Empirically the invoiced amount for month $M$ matches the bank payment recorded in
month $M\!+\!1$. The report therefore states program cost on the accrual basis,
with cash retained for the reserve position so reserves still tie to money on hand;
the cumulative difference is the outstanding payable, \ilsaProgramPayable{}. The
effective per-child program cost is the ratio of invoiced cost to children served
over billed months,
\begin{equation}
c \;=\; \frac{\sum_M \text{accrual}(M)}{\sum_M \text{kids}(M)} \;\approx\; \ilsaCohortCostPerChild{}\ \text{per child served},
\end{equation}
which blends the age-group books, welcome and graduation books, and per-piece
mailing.

\section{Trend estimation by ordinary least squares}

Forecastable series are extrapolated with a first-degree ordinary least-squares
(OLS) fit on the sequential month index. For observations $\{(t_i, y_i)\}_{i=0}^{n-1}$
with $t_i = i$, the line $\hat{y} = \beta_0 + \beta_1 t$ minimizes the residual sum
of squares $\sum_i (y_i - \beta_0 - \beta_1 t_i)^2$, giving the closed form
\begin{equation}
\beta_1 \;=\; \frac{\sum_i (t_i-\bar{t})(y_i-\bar{y})}{\sum_i (t_i-\bar{t})^2},
\qquad
\beta_0 \;=\; \bar{y} - \beta_1\,\bar{t}.
\end{equation}
Forward values are clamped to be non-negative, $\hat{y}_t = \max(\beta_0 + \beta_1 t,\,0)$.
This estimator is applied to three monthly series: recurring giving (slope
\ilsaRecurringSlope{}/month), invoiced program cost (slope \ilsaProgramSlope{}/month),
and the enrollment inflow \texttt{LETC} (slope \ilsaCohortLetcSlope{}/month). Because
\texttt{LETC} contains two campaign-driven spikes, its linear slope is upward-biased
relative to the non-campaign baseline; the corresponding projection is read as
upper-leaning, and the $\pm$\ilsaScenBand{}/year band brackets the uncertainty.

\section{Reserve projection and scenario bands}

Let $R_0$ be the current cash reserve and $\hat{I}_k,\hat{E}_k$ the OLS
extrapolations of recurring income and program cost for future month $k$. The
baseline reserve path accumulates projected net cash flow,
\begin{equation}
R_k \;=\; R_0 \;+\; \sum_{j=1}^{k} \bigl(\hat{I}_j - \hat{E}_j\bigr).
\end{equation}
The optimistic and pessimistic scenarios scale each month's net flow by a factor
that compounds with elapsed years $\tau_k = k/12$,
\begin{equation}
R_k^{\pm} \;=\; R_0 \;+\; \sum_{j=1}^{k} \bigl(\hat{I}_j - \hat{E}_j\bigr)\,(1 \pm \beta)^{\tau_j},
\qquad \beta = \ilsaScenBand{},
\end{equation}
where the pessimistic case uses $(1+\beta)$ (a widening deficit) and the optimistic
case $(1-\beta)$. No new grants are assumed, so all three paths are deliberately
conservative. The horizon runs to \ilsaScenEnd{}; under the baseline, reserves
reach \ilsaProjCashEnd{} at the horizon.

\section{The age-structured cohort model}

The forward enrollment and commitment figures come from a discrete-time,
age-structured population model --- a Leslie progression matrix without fecundity,
driven by an exogenous recruitment term. Enrollment is tracked in monthly
age-of-enrollment classes $a \in \{0,1,\dots,71\}$ (six annual groups
$\times$ twelve months); let $\mathbf{n}(t) \in \mathbb{R}^{72}$ be the population
vector, with group $g$ occupying classes $[12(g-1),\,12g)$. The model is seeded from
the latest group pyramid, distributed uniformly within each group,
$n_a(0) = G_g/12$ for $a$ in group $g$.

Each month every child ages one class, the oldest class graduates out, and new
children are recruited according to an intake distribution $\mathbf{w}$ (with
$\sum_a w_a = 1$):
\begin{equation}
\mathbf{n}(t+1) \;=\; \mathbf{A}\,\mathbf{n}(t) \;+\; R(t)\,\mathbf{w},
\qquad
\mathbf{A} =
\begin{pmatrix}
0 & 0 & \cdots & 0 & 0 \\
1 & 0 & \cdots & 0 & 0 \\
0 & 1 & \cdots & 0 & 0 \\
\vdots & & \ddots & & \vdots \\
0 & 0 & \cdots & 1 & 0
\end{pmatrix},
\label{eq:cohort}
\end{equation}
where the sub-diagonal shift matrix $\mathbf{A}$ advances each class and discards
class~71 (graduation). The number graduating in month $t$ is $n_{71}(t)$; total
enrollment is $N(t) = \sum_a n_a(t)$ and monthly program cost is $c\,N(t)$.

\paragraph{Recruitment.} $R(t)$ is the OLS extrapolation of the observed
\texttt{LETC} inflow (clamped at zero), optionally scaled by the
$(1\pm\beta)^{\tau}$ band. Graduation is endogenous --- it emerges from aging
Equation~\eqref{eq:cohort}, not from a separate assumption.

\paragraph{Intake age.} The \emph{age} at which new children enter is not directly
observed, so two intake distributions bracket it. The \emph{today's-mix} scenario
spreads recruits across classes in proportion to the current pyramid,
$w_a = p_g/12$ for $a$ in group $g$ (where $p_g$ is the current stock share of
group $g$); the \emph{birth-fed} scenario routes all recruits to birth,
$\mathbf{w} = \mathbf{e}_0$. Because the current pyramid is inverted (older-skewed),
today's mix graduates sooner and birth-fed carries a full six-year tail; the two
bracket the long-run liability.

\paragraph{Commitment runoff.} Setting $R \equiv 0$ runs the existing population off
to graduation and measures the locked-in obligation. A child in class $a$ owes
$72-a$ further monthly books, so the remaining commitment is
\begin{equation}
\text{Commitment} \;=\; c \sum_{a=0}^{71} (72 - a)\, n_a,
\end{equation}
which evaluates to \ilsaCohortRunoffTotal{} for the current roster --- the floor on
future program cost, assuming no new enrollment and no attrition.

\section{Reproducibility and limitations}

The pipeline is deterministic and parameterized by reporting period and projection
horizon; re-running the scripts and compiling twice reproduces the report exactly.
Sources are plain-text and diffable. The principal limitations are stated honestly:
linear extrapolation is sensitive to the two enrollment campaigns in the
\texttt{LETC} series; the intake age is unobserved (hence the two-scenario band);
attrition and voluntary disenrollment are not modeled; the recurring-giving trend is
sensitive to a small number of large early gifts; and the cash/accrual payable
reflects payment timing rather than any discrepancy in the underlying records.
